\section{Signed sums and denominator patterns}
\label{sec:denominator-patterns}

We refine the prime-power comparisons by retaining the signs of all
four radii at each denominator prime.  Use the primitive normalization
of Section~\ref{sec:introduction}:
\[
 (x,y)=(X/L,Y/L),\quad A_i=Lr_i,\quad
 G=\gcd(A_0,A_1,A_2,A_3),\quad H=L/G^2,\quad a_i=A_i/G.
\]
After exchanging coordinates, $X$ is odd, $4\mid Y$, $24\mid L$,
and the $A_i$ are odd.  Every prime of $G$ is $1$ modulo $4$,
and the $a_i$ are units at primes dividing $H$
(Proposition~\ref{int:content}).  Put
\[
 M_x=L-2X,\quad M_y=L-2Y,\qquad
 e_2=v_2(L)\geq3,\quad j=v_2(M_y)\geq3.
\]
The midline and diagonal exclusions make $M_x,M_y,X-Y,X+Y-L$
nonzero.  Also $v_2(M_x)=1$.

Let $\mathcal P=\{\pm1\}^4/\{\pm(1,1,1,1)\}$ and write
\[
 \alpha=(+,+,-,-),\quad \beta=(+,-,+,-),\quad
 \gamma=(+,-,-,+),\qquad
 \mathcal V=\{1,\alpha,\beta,\gamma\}.
\]
The elements of $\mathcal V$ have an even number of negative signs.
The other coset is represented by $o_i$, with a single negative
sign at corner $i$.  Define, up to an overall sign,
\[
 \Sigma_\epsilon=\sum_i\epsilon_i A_i,
 \qquad \sigma_\epsilon=\Sigma_\epsilon/G,
 \qquad P=A_0A_3-A_1A_2,quad Q=A_0A_3+A_1A_2.
\]
For the displayed representatives, direct expansion gives
\begin{equation}
 \begin{gathered}
 \Sigma_1\Sigma_\gamma=2P=-\Sigma_\alpha\Sigma_\beta,
 \qquad
 \Sigma_{o_0}\Sigma_{o_3}=2Q=\Sigma_{o_1}\Sigma_{o_2},\\
 PQ=-L^2M_xM_y,\\
 (Q^2-L^4)(L^4-P^2)
   =4L^4(X-Y)^2(X+Y-L)^2.
 \end{gathered}
 \label{eq:pat-product-identities}
\end{equation}

\begin{lemma}[Real bounds]
\label{lem:pat-real-bounds}
For a strict rational-distance point, all eight signed sums are nonzero,
and
\[
 \begin{aligned}
 2\sqrt2L&\leq\Sigma_1<(2+\sqrt2)L,\\
 |\Sigma_\alpha|,|\Sigma_\beta|&<\sqrt2L,
 &|\Sigma_\gamma|&<(2-\sqrt2)L,\\
 0<\Sigma_{o_i}&<(2+\sqrt2)L,
 &|P|&<L^2,\qquad L^2<Q<\frac{\sqrt5}{2}L^2.
 \end{aligned}
\]
\end{lemma}

\begin{proof}
The opposite-pair triangle inequalities give
$r_0+r_3\geq\sqrt2$ and $r_1+r_2\geq\sqrt2$.
Convexity of the sum of distances on the closed square gives
$r_0+r_1+r_2+r_3<2+\sqrt2$ in its interior.
Every adjacent-pair sum is strictly between $1$ and $1+\sqrt2$:
the lower bound is the triangle inequality and the upper bound
follows by convexity from its four vertex values.  In the
coordinates of \eqref{eq:sym-forward} this gives
$|h|,|k|<\sqrt2-1$.  With the present sign convention,
\[
 \Sigma_\alpha=-h\Sigma_1,\qquad
 \Sigma_\beta=-k\Sigma_1,\qquad
 \Sigma_\gamma=-hk\Sigma_1.
\]
These prove the stated even-sum bounds.  Each odd sum is positive:
for example $r_3<r_1+1<r_1+r_0+r_2$.
Its upper bound follows from that for the total sum.
Since $M_xM_y\ne0$, \eqref{eq:pat-product-identities} gives
$P\ne0$, and all the even sums are nonzero as well.

Put $\zeta=(x-\tfrac12)+i(y-\tfrac12)$, and write
$\zeta^2=A+iB$.  Then
\[
 \frac Q{L^2}=|\zeta^2-i/2|+|\zeta^2+i/2|\geq1.
\]
Equality lies on a diagonal, which is excluded.  Here is an
elementary upper bound.  Set
$u=(x-\tfrac12)^2$, $v=(y-\tfrac12)^2$, so
$0\leq u,v<1/4$, $A=u-v$ and $B^2=4uv$.
The separately convex polynomial
\[
 16(u-v)^2+\frac{64}{5}uv
\]
has values $0,1,1,4/5$ at the four vertices of
$[0,1/4]^2$, and is strictly less than $1$ in the stated range.
Thus $16A^2+(16/5)B^2<1$.  With
$R=A^2+B^2+1/4=(u+v)^2+1/4<1/2$, this is equivalent to
\[
 2R+2\sqrt{R^2-B^2}<\frac54.
\]
The left side is $(Q/L^2)^2$, proving the upper bound.
Finally, $PQ=-L^2M_xM_y$, $Q>L^2$ and $|M_x|,|M_y|<L$
give $|P|<L^2$.
\end{proof}

\begin{theorem}[Patterns and exact valuations]
\label{thm:pat-valuations}
For every odd prime $p\mid H$, put $e_p=v_p(H)$.  There is a
unique $\epsilon_p\in\mathcal P$ such that
\[
 a_i\equiv\epsilon_p(i)\epsilon_p(0)a_0\pmod {p^{e_p}}.
\]
For $p=2$, let $\epsilon_2(i)\epsilon_2(0)$ be the sign of
$A_i/A_0$ modulo $4$.  The same congruence then holds modulo
$2^{e_2}$.

For odd $p\mid H$, with $\epsilon=\epsilon_p$, the valuations are
\begin{equation}
\begin{aligned}
 v_p(\sigma_\epsilon)&=0,\\
 v_p(\sigma_{\epsilon\alpha})&=e_p+v_p(M_y),\\
 v_p(\sigma_{\epsilon\beta})&=e_p+v_p(M_x),\\
 v_p(\sigma_{\epsilon\gamma})&=2e_p+v_p(M_xM_y),\\
 v_p(\sigma_\eta)&=0\quad(\eta\notin\epsilon\mathcal V).
\end{aligned}
\label{eq:pat-odd-table}
\end{equation}
At $2$ the corresponding five values are, in the same order,
\begin{equation}
 2,\quad e_2+j,\quad e_2+1,\quad 2e_2+j-1,\quad 1.
 \label{eq:pat-two-table}
\end{equation}

Put $P_r=P/G^2$ and $Q_r=Q/G^2$.  If $\epsilon_p$ is even,
then for odd $p\mid H$
\[
 v_p(P_r)=2e_p+v_p(M_xM_y),\qquad v_p(Q_r)=0;
\]
if it is odd, the two values are exchanged.  At $2$ they are
$2e_2+j$ and $1$, respectively, with the same exchange rule.
For $p\mid G$ with $p\nmid H$, both $P_r$ and $Q_r$ are units.
\end{theorem}

\begin{proof}
All six squared differences $a_i^2-a_j^2$ are divisible by $H$.
For odd $p$, their unit roots therefore differ by unique signs
modulo $p^{e_p}$.  At $2$, choose signs so that the four resulting
$b_i$ are congruent to $A_0$ modulo $4$.  Then $v_2(b_i+b_j)=1$.
The horizontal squared differences have valuation $e_2+1$,
and the vertical ones have valuation $e_2+j$; hence the four
$b_i$ are congruent modulo $2^{e_2}$.

For odd $p$, put $b_i=\epsilon_p(i)\epsilon_p(0)a_i$.
All $b_i$ are congruent to $a_0$ modulo $p^{e_p}$.
The sum $S'=b_0+b_1+b_2+b_3$ and
$Q'=b_0b_3+b_1b_2$ are units.  With
$P'=b_0b_3-b_1b_2$ and $T'=b_0-b_1-b_2+b_3$, one has
\[
 S'T'=2P',\qquad P'Q'=-H^2M_xM_y.
\]
This proves the valuation of $\sigma_{\epsilon\gamma}$.
Moreover,
\[
 (b_2-b_0)+(b_3-b_1)
 =H M_y\left(\frac1{b_2+b_0}+\frac1{b_3+b_1}\right).
\]
The bracket is a unit modulo an odd $p$, since it is congruent
to $a_0^{-1}$.  This gives the $\alpha$ entry; the horizontal
version gives the $\beta$ entry.  The other coset consists,
up to sign, of $S'-2b_i$, which are units.

At $2$, the same reciprocal bracket has valuation zero: after
factoring out $1/2$, its two unit terms agree modulo $4$.
Thus the $\alpha$ and $\beta$ entries are $e_2+j$ and
$e_2+1$.  The total sum has valuation two, and the other coset
has valuation one.  The product identities then give
$v_2(P')=2e_2+j$ and $v_2(T')=2e_2+j-1$.

The normalized $P'$ is $\pm P_r$ for an even pattern and
$\pm Q_r$ for an odd one, giving the exchange rule.  Finally,
if $p\mid G$ and $p\nmid H$, then $p$ is split and $X,Y$,
$M_x,M_y$ are units at $p$.  Since
$P_rQ_r=-H^2M_xM_y$, both factors are units.
\end{proof}

\begin{proposition}[Colour determines the nonsplit patterns]
\label{prop:pat-colours}
For an odd prime $p\mid L$ with $p\equiv3\pmod4$,
\[
 \epsilon_p(i)\epsilon_p(j)=\leg{m_im_j}{p}.
\]
For $p\mid L$ with $p\equiv1\pmod4$, all six symbols
$\leg{m_im_j}{p}$ are $1$.  At $2$,
$\epsilon_2(i)\epsilon_2(j)=1$ if and only if
$m_i\equiv m_j\pmod4$.  Consequently, for nonsplit odd $p$
the pattern is even exactly when $\leg np=1$, and
$\epsilon_2$ is even exactly when $n\equiv1\pmod8$.
\end{proposition}

\begin{proof}
Use $U_i=m_iq_i^2$ and $A_i=m_i\Nm(q_i)$ from the Gaussian
normalization.  At an inert prime, the four $U_i$ have the same
nonzero image in $\F_{p^2}$.  If $\mu_i^2=m_i$, then
$\mu_iq_i=\pm\mu_0q_0$.  Frobenius sends
$\mu_i$ to $\leg{m_i}{p}\mu_i$, so taking norms gives
$\leg{m_i}{p}A_i=\leg{m_0}{p}A_0$ modulo $p$.
The sign lift is unique, proving the first assertion.
At a split prime choose the factor of $p$ at which $U_0$ is a
unit.  Reduction in its residue field $\F_p$ gives
$m_i/m_j=(q_j/q_i)^2$.
At $2$, each odd Gaussian norm $\Nm(q_i)$ is $1$ modulo $4$,
so $A_i\equiv m_i\pmod4$.
Taking the product of the four signs proves the parity claims,
using $n\equiv1$ or $7\pmod8$.
\end{proof}

Let $E_+$ be the product of $p^{e_p}$ over odd primes of $H$
with even pattern, and $E_-$ the corresponding product for odd
patterns.  Thus $H=2^{e_2}E_+E_-$.

\begin{theorem}[The two product denominators]
\label{thm:pat-product-denominators}
Write $t=r_0r_3+r_1r_2=N_t/D_t$ and
$z=r_0r_3-r_1r_2=N_z/D_z$ in lowest terms with positive
denominators.  If $n\equiv1\pmod8$, then
\[
 D_t=(2^{e_2}GE_+)^2/2,\qquad D_z=(GE_-)^2;
\]
if $n\equiv7\pmod8$, then
\[
 D_t=(GE_+)^2,\qquad D_z=(2^{e_2}GE_-)^2/2.
\]
In both cases $D_tD_z=L^2/2$ and $\gcd(D_t,D_z)=G^2$.

There are factorizations $n=n_3n_4=n_1n_2$ and positive odd
integers $u,v,u',v'$ such that
\begin{equation}
 2D_t=n_4v^2-n_3u^2,\qquad
 2D_z=n_1u'^2+n_2v'^2,
 \label{eq:pat-norm-denominators}
\end{equation}
and
\[
 n u v u'v'=|(X-Y)(X+Y-L)|,\qquad
 0<\frac{n_3u^2}{n_4v^2}<9-4\sqrt5.
\]
These two factorizations of $n$ are not required to be the
corner-colour factorizations.
\end{theorem}

\begin{proof}
We have $t=Q_r/(G^2H^2)$ and $z=P_r/(G^2H^2)$.
The valuation table cancels precisely the square factors belonging
to the opposite pattern parity.  At a prime shared by $G$ and $H$,
$M_xM_y$ is a unit, so the residual denominator contains exactly
the asserted factor from $G^2$.  At a prime of $G$ outside $H$
both numerators are units.  The dyadic entries are $1$ and
$2e_2+j$, and Proposition~\ref{prop:pat-colours} identifies
which is which.  This proves the denominator formulas, including
their product and gcd.  The numerator opposite an odd denominator
is even; the numerator opposite an even denominator is odd.

The kernel identity gives $t^2-1\in n\Q^{\times2}$.
The last identity of \eqref{eq:pat-product-identities} gives
$1-z^2\in n\Q^{\times2}$.  Both are positive by
Lemma~\ref{lem:pat-real-bounds}.  The positive odd integers
$N_t-D_t,N_t+D_t$ are coprime, as are
$D_z-N_z,D_z+N_z$.  Their products have squarefree part $n$,
so write
\[
 N_t-D_t=n_3u^2,\quad N_t+D_t=n_4v^2,\qquad
 D_z-N_z=n_1u'^2,\quad D_z+N_z=n_2v'^2.
\]
This proves \eqref{eq:pat-norm-denominators}.
Multiplying the four equations and using $D_tD_z=L^2/2$
proves the product formula for $uv u'v'$.
Finally $(t-1)/(t+1)<(\sqrt5-2)/(\sqrt5+2)=9-4\sqrt5$.
\end{proof}

\begin{corollary}[At least three denominator patterns]
\label{cor:pat-three-patterns}
If $G=1$, the primes dividing $L$ realize at least three different
patterns, and both cosets of $\mathcal V$ occur.  In particular,
if only one corner colour is nontrivial, then $L$ has a prime
divisor congruent to $1$ modulo $4$.
\end{corollary}

\begin{proof}
Put $m_p=v_p(M_xM_y)$ for odd $p\mid H$.
If all patterns are even, the table gives
\[
 |P_r|\geq 2^jH^2\prod_{p\mid H,\ p\ {
m odd}}p^{m_p},
 \qquad |P_r|<G^2H^2.
\]
If all are odd, use $Q_r$ and its upper bound
$(\sqrt5/2)G^2H^2$ instead.  Since $2^j\geq8$, $3\mid H$
and $m_3\geq1$, either alternative is impossible for $G=1$.

Suppose there are only two patterns $\epsilon,\epsilon'$ in
different cosets.  Apply the valuation table to
$\sigma_{\epsilon\gamma}\sigma_{\epsilon'\gamma}$.
Each odd prime contributes $2e_p+m_p$, and the dyadic contribution
is $2e_2+j$.  Therefore
\[
 2^jH^2\prod_{p\mid H,\ p\ {
m odd}}p^{m_p}
 \leq|\sigma_{\epsilon\gamma}\sigma_{\epsilon'\gamma}|
 <(2+\sqrt2)^2G^2H^2.
\]
For $G=1$ this contradicts $24>(2+\sqrt2)^2$.

For the last assertion, if $G>1$ a prime of $G$ suffices.
If $G=1$ and no split prime divides $L$, all patterns are
determined by the colours.  When only $m_k$ is nontrivial they
belong to $\{1,o_k\}$, contrary to the first assertion.
\end{proof}

\begin{corollary}[A sharper dyadic comparison]
\label{cor:pat-dyadic-bound}
If $m_i\equiv1\pmod4$ for every corner, and
$K=2^{v_2(L)}$, then
\[
 K^2<\frac{2-\sqrt2}{4}\,L.
\]
\end{corollary}

\begin{proof}
Here $\epsilon_2=1$.  The exact valuation of $\Sigma_\gamma$
is $2e_2+j-1\geq2e_2+2$, so
$4K^2\leq|\Sigma_\gamma|<(2-\sqrt2)L$.
\end{proof}
